With one complete recipe execution, the 50-feature network gave 0.19\% mean improvement over the best fixed training recipe, with a 95\% interval of [-0.40\%, 0.79\%]. The 24-feature network gave 0.44\% with interval [-0.28\%, 1.10\%]. Both point estimates are close to the fixed-recipe result; the intervals include zero, leaving the direction of the average difference uncertain. For the 50-feature network, mean area fell by 0.31\% and mean delay rose by 0.47\% relative to the fixed recipe.

Running the full network's five highest-ranked recipes gave 5.12\% improvement over the fixed recipe and 3.99\% mean oracle regret. It reached the exact library optimum on 15.28\% of held-out predictions and came within 5\% of it on 70.24\%. These rates average the three held-out predictions per circuit. The corresponding mean area and delay reductions were 2.44\% and 7.17\%.

The global top-five baseline improved the fixed-recipe objective by 5.01\% with 4.13\% regret. The NN top-five result improved this global shortlist by 0.06\% on average (95\% interval [-0.36\%, 0.48\%]). This interval also includes zero. The experiments show how a small candidate budget recovers much of the available quality; the global shortlist remains a strong inexpensive alternative.
